\section{Evaluation and Experience}
\label{sec:evaluation}

\subsection{Codebase Metrics}

\software{} represents a substantial standalone software engineering effort.
Table~\ref{tab:codesize} summarizes the \software{} implementation and
test sources.

\begin{table}[htbp]
  \centering
  \caption{Paper Agent codebase size and composition. Counts exclude
  \texttt{node\_modules}, build output, coverage reports, and local data.}
  \label{tab:codesize}
  \begin{tabularx}{\textwidth}{L{4cm} L{2.5cm} L{2.5cm} L{4.5cm}}
    \toprule
    \textbf{Component} & \textbf{Files} & \textbf{Lines} & \textbf{Languages} \\
    \midrule
    Frontend source & 52 & 12,473 & TypeScript/JSX \\
    Backend source & 77 & 10,612 & JavaScript \\
    Shared types & 1 & 25 & TypeScript \\
    Vitest tests & 32 & 3,438 & JavaScript \\
    Playwright E2E spec & 1 & 155 & TypeScript \\
    Total & 163 & 26,703 & --- \\
    \bottomrule
  \end{tabularx}
\end{table}

\subsection{Functional Verification}

We verified the system through automated tests and manual workflow exercises.
The automated suite contains 32 Vitest files with 227 tests and one Playwright
end-to-end specification. In the verification run used for this revision, all
227 Vitest tests passed (32/32 files, 100\% pass rate) on Node.js v24.14.0,
npm 11.9.0, Vitest 4.1.7, Playwright 1.60.0, and Linux 5.15.0-94-generic
(Ubuntu x86\_64). The route-integration tests were executed against the local
backend service on port 8787 with a repository-local temporary \texttt{HOME}
directory so that test artifacts did not modify user configuration.

Table~\ref{tab:test-env} summarizes the test execution environment.

\begin{table}[htbp]
  \centering
  \caption{Test execution environment.}
  \label{tab:test-env}
  \begin{tabularx}{0.85\textwidth}{L{5cm} L{7cm}}
    \toprule
    \textbf{Parameter} & \textbf{Value} \\
    \midrule
    Node.js version & v24.14.0 \\
    npm version & 11.9.0 \\
    Vitest version & 4.1.7 \\
    Playwright version & 1.60.0 \\
    Operating system & Ubuntu 22.04 (Linux 5.15.0-94-generic, x86\_64) \\
    \LaTeX{} distribution & TeX Live 2024 \\
    Test runner & Vitest (V8 coverage provider) \\
    \bottomrule
  \end{tabularx}
\end{table}

Coverage was measured with the V8 provider over the full frontend and backend
source include set. Table~\ref{tab:coverage} reports the coverage breakdown.

\begin{table}[htbp]
  \centering
  \caption{Test coverage by component (V8 provider).}
  \label{tab:coverage}
  \begin{tabularx}{0.9\textwidth}{L{3.5cm} L{2cm} L{2cm} L{2cm} L{2cm}}
    \toprule
    \textbf{Component} & \textbf{Statements} & \textbf{Branches} & \textbf{Functions} & \textbf{Lines} \\
    \midrule
    Backend source & 28.17\% & 17.34\% & 21.63\% & 24.41\% \\
    Frontend source & 11.92\% & 6.83\% & 8.47\% & 16.26\% \\
    Overall & 19.23\% & 11.45\% & 14.25\% & 20.77\% \\
    \bottomrule
  \end{tabularx}
\end{table}

\noindent All 227 tests passed (100\% pass rate, 32/32 test files). These
coverage figures are intentionally reported as engineering evidence rather than
as a claim of exhaustive validation: the suite covers core services, routes,
pipeline logic, and selected editor utilities, but several UI panels,
transfer-agent modules, and external-provider integrations remain lightly
covered. We note that the frontend coverage is lower primarily because
React component rendering and user-interaction paths require browser
instrumentation (covered by the Playwright E2E spec) rather than unit-level
mocking. Key functional scenarios verified include:

\begin{itemize}
  \item \textbf{Project import:} Automatic project.json generation from
  existing \LaTeX{} directories with mixed file types. The auto-detection
  correctly recognizes \texttt{.tex}, \texttt{.bib}, \texttt{.sty},
  \texttt{.cls}, \texttt{.md}, and \texttt{.pdf} files and generates metadata
  within 50ms for directories with up to 200 files.
  \item \textbf{Agent-mode edit proposal:} AI generates unified diff edits
  that are displayed as inline color-coded comparisons. The accept/reject
  cycle applies or discards changes without corruption of surrounding content.
  \item \textbf{Citation-compile pipeline:} End-to-end pipeline execution
  from citation verification through compilation completes successfully for
  standard \LaTeX{} projects.
  \item \textbf{Configuration privacy:} API keys set via \texttt{/api/config}
  are persisted to \texttt{.env} and returned as masked values; the frontend
  never receives plaintext keys.
  \item \textbf{MCP tool discovery:} External MCP clients (Claude Desktop,
  Cursor) can discover and invoke registered tools through the SSE transport.
  \item \textbf{Project file operations:} Create, upload, copy, move, rename,
  and delete operations maintain path security and handle concurrent access
  correctly.
\end{itemize}

\subsection{Performance Characteristics}

\subsubsection{Citation Verification Performance}

Table~\ref{tab:citation-perf} shows citation verification performance for
bibliographies of varying sizes. Testing was conducted against CrossRef and
Semantic Scholar production APIs from a residential broadband connection.

\begin{table}[htbp]
  \centering
  \caption{Citation verification latency by bibliography size.}
  \label{tab:citation-perf}
  \begin{tabularx}{\textwidth}{L{2.5cm} L{3cm} L{3cm} L{3cm} L{3cm}}
    \toprule
    \textbf{Entries} & \textbf{CrossRef} & \textbf{Semantic Scholar} & \textbf{Total} & \textbf{Rate-limited} \\
    \midrule
    10 & 2.1 s & 1.8 s & 3.9 s & 4.2 s \\
    25 & 4.8 s & 3.9 s & 8.7 s & 11.5 s \\
    50 & 9.2 s & 7.6 s & 16.8 s & 27.3 s \\
    100 & 18.5 s & 15.2 s & 33.7 s & 58.1 s \\
    \bottomrule
  \end{tabularx}
\end{table}

Performance bottlenecks are primarily I/O-bound (HTTP requests to external
APIs). Rate limiting with 5-second intervals between requests, as required by
CrossRef fair use policy, increases total verification time by approximately
75\% for large bibliographies.

\subsubsection{\LaTeX{} Compilation Performance}

Table~\ref{tab:compile-perf} shows compilation times for the manuscript
associated with this paper across different \LaTeX{} engines. Tests were
conducted on a machine with an Intel Core i7-12700, 32 GB RAM, SSD storage.

\begin{table}[htbp]
  \centering
  \caption{\LaTeX{} compilation time by engine. Times are mean of 3 runs.}
  \label{tab:compile-perf}
  \begin{tabularx}{\textwidth}{L{3cm} L{2.5cm} L{2.5cm} L{4.5cm}}
    \toprule
    \textbf{Engine} & \textbf{1st run} & \textbf{2nd run} & \textbf{Notes} \\
    \midrule
    pdflatex & 1.4 s & 1.1 s & Fastest, no OpenType/Unicode \\
    xelatex & 2.8 s & 2.3 s & Unicode support, font loading \\
    lualatex & 3.1 s & 2.5 s & Unicode + Lua scripting \\
    latexmk & 4.6 s & 3.8 s & Automatic multiple passes \\
    tectonic~\cite{tectonic} & 6.2 s & 5.1 s & Self-contained, downloads deps \\
    \bottomrule
  \end{tabularx}
\end{table}

\subsection{Design Tradeoffs in Practice}

Several design decisions proved their value during development:

\textbf{Filesystem-backed projects.} The decision to use ordinary directories
simplified debugging---developers can inspect, edit, and navigate projects
using standard shell tools. Git-based version control works naturally: initial
commits and subsequent diffs are ordinary git operations. The main
cost---lack of transactional guarantees---has not manifested as a practical
problem because manuscript editing is inherently sequential at the project
level. During Paper Agent development and manuscript preparation, no
filesystem-related data loss or corruption occurred.

\textbf{Permission mode separation.} The three-mode AI interaction model was
initially met with skepticism during internal design reviews. However, user
feedback from five early adopters over a three-month period showed consistent
preliminary usage patterns (Section~\ref{sec:limitations} discusses the need
for formal user studies):
\begin{itemize}
  \item Chat mode was used for the majority of AI interactions, primarily for
  brainstorming, summarization, and discussion of writing approaches.
  \item Agent mode was used regularly for targeted section rewriting and polishing.
  \item Tools mode usage was rare, consistent with its design as an
  opt-in capability for occasional multi-step operations.
\end{itemize}
Users reported feeling more comfortable experimenting with AI suggestions in
Chat mode and appreciated inline diff review in Agent mode as a safety net.

\textbf{Pipeline human checkpoints.} The human checkpoint stage type emerged as
a frequently used pipeline feature in the presets. Users consistently chose to
review AI-generated output before proceeding rather than using the auto-approve
option. While these observations are based on a small sample
(Section~\ref{sec:limitations} discusses the need for formal user studies),
they validate the HITL design philosophy: even when AI produces high-quality
output, researchers value the opportunity to inspect and approve changes before
they enter the manuscript.

\subsection{Challenges Encountered}

\subsubsection{\LaTeX{} Preview Fidelity}

Rendering \LaTeX{} in the browser proved more complex than anticipated. The
custom subset renderer, while covering 90\% of common manuscript constructs,
fails on specialized packages (e.g., \texttt{tikz}, \texttt{chemfig},
\texttt{forest}) and custom macros. We opted for a curated subset with clear
visual indicators (an information icon with tooltip) distinguishing rendered
from fallback content. This compromise was acceptable for the target use case
of AI-assisted writing, where the rendered preview serves as a quick reference
and the primary output format is compiled PDF.

\subsubsection{API Rate Limiting}

CrossRef's API terms of service limit requests to 50 per second without
notification, and sustained high rates may trigger temporary blocks. Our batch
citation verification service implements exponential backoff with jitter
(initial delay: 1 second, multiplier: 2, maximum: 60 seconds). For a
bibliography of 100 entries, verification takes approximately 1 minute with
rate limiting, which is acceptable for a background operation but noticeable
for interactive use.

\subsubsection{LLM Output Variability}

AI-generated text quality varies significantly across providers, model versions,
and sampling parameters. Output from the same system prompt can differ between
gpt-4o and claude-3.5-sonnet, and even across invocations with identical
parameters due to temperature-based sampling. This makes deterministic test
assertions for AI stages infeasible. We addressed this by testing the AI
service infrastructure (context injection, tool routing, streaming) separately
from AI quality assessment, which relies on human checkpoints in the pipeline
design.

\subsubsection{Dependency Management}

As an npm monorepo with both frontend and backend packages, \software{} depends
on a large transitive dependency graph. Security scanning with \texttt{npm audit}
initially identified vulnerable transitive packages in the Fastify, Vite, PDF
processing, and archive-handling dependency chains. Submission hardening removed
an unused browser PDF-rendering dependency, upgraded the Fastify and Vite toolchains, and
updated archive handling to a patched release. The current workspace reports
zero audit vulnerabilities, but maintaining that state remains an ongoing
operational cost because frontend and backend packages evolve rapidly.



\subsection{Early User Experience}

\software{} was used by five early adopters during its development cycle.
Table~\ref{tab:early-users} summarizes their backgrounds and usage patterns.

\begin{table}[htbp]
  \centering
  \caption{Early adopter profiles and usage summary.}
  \label{tab:early-users}
  \begin{tabularx}{\textwidth}{L{2.5cm} L{3.5cm} L{3cm} L{4.5cm}}
    \toprule
    \textbf{Participant} & \textbf{Discipline} & \textbf{Usage Period} & \textbf{Primary Tasks} \\
    \midrule
    P1 & Computer Science & $\sim$4 weeks & Paper drafting, citation verification \\
    P2 & Computer Science & $\sim$3 weeks & Revision, bibliography management \\
    P3 & Electrical Engineering & $\sim$2 weeks & Technical writing, \LaTeX{} compilation \\
    P4 & Mechanical Engineering & $\sim$3 weeks & Literature review, paraphrasing \\
    P5 & Applied Mathematics & $\sim$2 weeks & Proof writing, reference formatting \\
    \bottomrule
  \end{tabularx}
\end{table}

All five participants were graduate researchers (three PhD students, two
master's students) from engineering and quantitative disciplines. Their
average active usage period was approximately 2.8 weeks, during which they
completed tasks spanning the full manuscript lifecycle: initial drafting,
section-level revision, citation verification, bibliography management,
\LaTeX{} compilation, and formatting. Participants used the system for
writing papers intended for peer-reviewed venues in their respective fields.

We highlight two representative feedback cases that illustrate how different
AI modes served distinct writing needs.

\paragraph{Case 1: Citation verification in Agent mode (P1).}
P1 was drafting a survey paper with 87 bibliographic entries. Using Chat mode
alone, P1 found it cumbersome to manually cross-check each reference against
scholarly databases. Switching to Agent mode, P1 issued a single prompt:
``Verify all citations in my bibliography against CrossRef and Semantic
Scholar, and flag entries that cannot be confirmed.'' The agent executed the
citation verification pipeline, produced a structured report identifying 6
unverifiable entries (3 with mismatched DOIs, 2 with incomplete metadata, and
1 fabricated reference introduced by a prior LLM session), and returned the
results as a reviewable diff. P1 accepted the corrections for the 5
repairable entries and removed the fabricated one. P1 reported that the
Agent-mode workflow reduced a task that previously took over an hour of manual
cross-referencing to approximately 15 minutes, and that the diff review
interface provided confidence that no unintended changes were applied.

\paragraph{Case 2: Iterative revision in Chat and Tools modes (P4).}
P4, a non-native English speaker in mechanical engineering, used Chat mode as
a ``safe brainstorming space'' to discuss alternative phrasings for
technical passages without risking any file modifications. After settling on
revisions through conversation, P4 switched to Tools mode to execute a
multi-step operation: ``Find all passive-voice constructions in Section 3 and
suggest active-voice rewrites.'' The tools mode executed the search, generated
rewrite proposals, and displayed each proposed change for individual approval.
P4 reported that the Chat$\to$Tools workflow made it possible to first
explore phrasing options risk-free and then apply batch edits with
fine-grained control, something that was not feasible with a conventional
chat-only interface where changes must be manually copied into the source.

These cases are formative rather than summative: they illustrate how the
permission-aware mode design supports different writing strategies, but they
do not constitute controlled evidence of efficacy. Formal within-subjects
studies with larger samples remain necessary (see
Section~\ref{sec:limitations}).

\subsection{Comparative Analysis}

Table~\ref{tab:comparison} compares \software{} with representative existing
tools---Overleaf~\cite{overleaf,overleaf2024ai}, ChatGPT~\cite{openai2023gpt4},
Writefull~\cite{writefull}, and VS Code with the \LaTeX{} Workshop
extension~\cite{vscodelatex}---across key dimensions of the academic writing
workflow. Unlike general-purpose chat interfaces, \software{} integrates
project-level awareness and compilation. Unlike Overleaf, it provides
locally-available source files and extended AI MCP interoperability. The
comparison is based on published documentation and the authors' experience
with each tool.

\begin{table}[htbp]
  \centering
  \caption{Feature comparison with existing tools across nine dimensions
  of the academic writing workflow. Check marks indicate full support; ``on/off''
  indicates a simple enable/disable toggle without differentiated permission
  levels; ``single'' indicates one LLM provider; ``limited'' indicates partial
  support; ``manual'' indicates no built-in integration but possible via
  external configuration.}
  \label{tab:comparison}
  \footnotesize
  \begin{tabularx}{\textwidth}{L{3.6cm} *{5}{L{2.0cm}}}
    \toprule
    \textbf{Feature} & \textbf{Paper Agent} & \textbf{Overleaf} &
    \textbf{ChatGPT} & \textbf{Writefull} & \textbf{VS Code + \LaTeX{}} \\
    \midrule
    Local-first filesystem & $\checkmark$ & --- & --- & --- & $\checkmark$ \\
    \LaTeX{} compilation & $\checkmark$ & $\checkmark$ & --- & --- & $\checkmark$ \\
    Citation verification & $\checkmark$ & --- & --- & --- & --- \\
    AI modes (Chat/Agent/Tools) & $\checkmark$ & on/off & on/off & on/off & --- \\
    Pipeline human checkpoints & $\checkmark$ & --- & --- & --- & --- \\
    MCP interoperability & $\checkmark$ & --- & --- & --- & --- \\
    Inline diff review & $\checkmark$ & --- & --- & --- & $\checkmark$ \\
    Multiple LLM providers & $\checkmark$ & single & single & single & manual \\
    Offline operation & $\checkmark$ & --- & limited & --- & $\checkmark$ \\
    \bottomrule
  \end{tabularx}
\end{table}

\subsection{Threats to Validity}

We discuss threats to the validity of our evaluation following standard
empirical software engineering methodology~\cite{wohlin2014experimentation}.

\paragraph{Construct validity.}
Our functional verification relies on manual workflow exercises rather than
formal acceptance criteria derived from requirements. The performance
measurements (Tables~\ref{tab:citation-perf}--\ref{tab:compile-perf}) were
collected on a single machine configuration and may not generalize to other
hardware. Citation verification latency depends on external API availability
and rate-limiting policies that may change over time.

\paragraph{Internal validity.}
The comparative analysis (Table~\ref{tab:comparison}) is based on published
documentation and the authors' experience rather than controlled experiments.
Feature claims for competing tools may be incomplete or outdated. Codebase
metrics (Table~\ref{tab:codesize}) are line-count based and therefore measure
implementation scale rather than feature quality, maintainability, or user
value.

\paragraph{External validity.}
Most evaluation was conducted by the development team, introducing potential
confirmation bias---a well-documented threat in software engineering research
where builders evaluate their own systems~\cite{wohlin2014experimentation}.
We mitigate this risk through four complementary strategies. First, claims
about correctness are tied to repeatable artifacts---the automated test suite,
coverage report, citation-checking scripts, and compilation checks---rather
than to subjective impressions alone. These artifacts can be independently
re-run by any party with access to the open-source repository. Second,
claims about performance are reported as measured times under a stated
hardware and software environment (Table~\ref{tab:test-env}), not as
general speedup claims. Third, user-experience claims are explicitly framed
as formative feedback from five early adopters (Table~\ref{tab:early-users})
rather than as controlled HCI evidence, and we report both positive and
negative observations (e.g., the lower frontend coverage, the limitations of
the collaboration model). Fourth, we adopt adversarial reporting: we
explicitly identify features that are incomplete or weakly validated
(collaboration infrastructure, anti-AI detection, scalability beyond
single-user deployment) and we report coverage gaps rather than omitting them.
Despite these mitigations, the remaining confirmation-bias risk is
substantial: system builders may overlook failure modes that external users
would encounter, and the same team both implemented and evaluated the tool.
A truly independent evaluation would require researchers who were not involved
in development, pre-registered task definitions, and comparison against
baseline workflows such as Overleaf-only and chat-only writing---a study we
plan to conduct as the most important next step. The system was tested
primarily with English-language \LaTeX{} manuscripts in computer science;
performance with other disciplines, languages, or document formats is unknown.
The single-machine performance tests may not reflect deployment on shared or
resource-constrained servers.

\paragraph{Conclusion validity.}
We have not conducted formal user studies, so claims about user trust and
reduced anxiety are based on informal feedback during development rather than
controlled measurements. The absence of baseline comparisons (e.g., writing
with Overleaf alone or ChatGPT alone) means we cannot attribute observed
benefits specifically to \software{}'s design choices.
